% Physical page 1 of the concise study, second half: exactly four overview
% rows, immediately after the map of the appointed elements. They orient a
% reader who has not yet met the three interpretations. Each interpretation's
% own four senses are developed in the expansive study, and the places where
% they differ are compared in the commentary that begins on the fifth page.
% Every parenthetical name stands beside what that author says of his own
% passage; where the passage is not the one the row speaks of, it is named.
% A clause that rests on one of two alternatives, the longer form of the
% Gospel or one of the two Communion antiphons, is marked with it.

\vspace{0.3em}
\zlabel{triptych:concise:overview:start}
\begin{conciseoverview}
Literal\zlabel{triptych:concise:sense:literal} & Isaiah promises a feast for
all peoples on the Lord's mountain, where death is cast down. Jesus tells the
chief priests and elders of a king whose invited refuse his son's wedding, one
to his farm and another to his merchandise, while the rest kill his servants;
the hall is filled from the roads, and in the longer form a guest without a
wedding garment is cast out. Paul has learned to be full and to be hungry.\\
Allegorical\zlabel{triptych:concise:sense:allegorical} & The king is the
Father, and the wedding the Church joined to the Son through the mystery of
the Incarnation (Gregory), gathered from Jews and gentiles alike (Jerome).
Aquinas sets Isaiah's feast on the mountain where Christ suffered beside the
king's ``my dinner is prepared''. In the longer form the garment is charity
(Augustine, Gregory), Christ put on (Aquinas), or the Spirit's glory (Hilary).\\
Moral\zlabel{triptych:concise:sense:moral} & Set spiritual things above even
necessary ones (Chrysostom); no temporal thing should keep anyone from God
(Aquinas), who on Philippians bids plenty and want be used alike with virtue.
In the longer form: let charity be born and grow, and clothe the naked
(Augustine); keep the garment the call gave (Chrysostom); be humble, not
knowing whether you are chosen (Gregory).\\
Anagogical\zlabel{triptych:concise:sense:anagogical} & The wedding is the
glory received in the resurrection (Hilary), when death is overcome and
Christ's kingdom comes (Jerome, on Isaiah), and all desire is filled (Aquinas,
on Philippians 4:19). In the longer form the king's coming in is the judgment
(Jerome, Gregory). Where the second Communion antiphon is chosen: the more
charity, the fuller the vision of him as he is (Aquinas).\\
\end{conciseoverview}
\zlabel{triptych:concise:overview:end}
